\chapter{Confidence Intervals}
\label{ch:confidence-intervals}

A dashboard says the mean page-load time this week is 212 ms. Is the true
mean 212.0? Almost certainly not. It is 212 plus or minus something, and the
size of that something decides whether last week's 205 ms is a real
regression or noise. Chapter~6 gave us the raw ingredient: the standard error
$\SE$, the typical distance between an estimate and the truth. This chapter
turns $\SE$ into a range with a stated success rate.

\section{From standard error to interval}
\label{sec:ci-from-se}

Take $n=36$ page loads. Their mean is $\xbar = 212$ ms. Suppose years of
monitoring tell us the per-request standard deviation is $\sigma = 30$ ms.
Then
\[
  \SE(\Xbar) = \frac{\sigma}{\sqrt n} = \frac{30}{\sqrt{36}} = \frac{30}{6} = 5 \text{ ms}.
\]
By the Central Limit Theorem (Chapter~6), $\Xbar$ is approximately Normal,
centred on the true mean $\mu$ with standard deviation 5 ms. A Normal variable
lands within $1.96$ standard deviations of its centre 95\% of the time. So
95\% of the time,
\[
  \abs{\Xbar - \mu} \le 1.96 \times 5 = 9.8 \text{ ms}.
\]
If $\Xbar$ is within 9.8 ms of $\mu$, then $\mu$ is within 9.8 ms of $\Xbar$.
That gives the interval
\[
  212 \pm 9.8 = [202.2,\ 221.8] \text{ ms}.
\]
Last week's 205 ms is inside it. This week's data are consistent with no
change at all.

\begin{figure}[H]
\centering
\begin{tikzpicture}
\begin{axis}[
  width=12cm, height=5.6cm,
  domain=-3.6:3.6, samples=121,
  axis lines=left, ymin=0, ymax=0.47,
  ytick=\empty,
  xtick={-3,-1.96,-1,0,1,1.96,3},
  xticklabels={$\mu-15$,$\mu-9.8$,$\mu-5$,$\mu$,$\mu+5$,$\mu+9.8$,$\mu+15$},
  xticklabel style={font=\small},
  xlabel={value of $\Xbar$ (ms)}, ylabel={density},
  clip=false,
]
\addplot[fill=histfill!60, draw=none, domain=-1.96:1.96] {0.39894*exp(-x^2/2)} \closedcycle;
\addplot[fill=trapred!35, draw=none, domain=-3.6:-1.96] {0.39894*exp(-x^2/2)} \closedcycle;
\addplot[fill=trapred!35, draw=none, domain=1.96:3.6] {0.39894*exp(-x^2/2)} \closedcycle;
\addplot[distblue, thick] {0.39894*exp(-x^2/2)};
\node at (axis cs:0,0.16) {\small 95\%};
\node[trapred] at (axis cs:-2.75,0.08) {\small 2.5\%};
\node[trapred] at (axis cs:2.75,0.08) {\small 2.5\%};
\draw[<->, thick] (axis cs:-1.96,0.43) -- (axis cs:1.96,0.43)
  node[midway, above, font=\small] {$\pm 1.96\,\SE = \pm 9.8$ ms};
\end{axis}
\end{tikzpicture}
\caption{The sampling distribution of $\Xbar$ for $n=36$, $\sigma=30$
($\SE=5$ ms). The middle 95\% of sample means fall within 9.8 ms of $\mu$.
That fact, read backwards, is the confidence interval.}
\label{fig:ci-sampling-shaded}
\end{figure}

Every interval in this chapter has the same skeleton.

\begin{definition}[Confidence interval, general form]
A \emph{confidence interval} for a parameter $\theta$ at confidence level
$1-\alpha$ is a range computed from the data,
\[
  \text{estimate} \;\pm\; (\text{critical value}) \times \SE(\text{estimate}),
\]
built so that, before the data are drawn, the random interval contains
$\theta$ with probability $1-\alpha$. The number $\alpha$ is the miss rate we
accept: $\alpha=0.05$ gives a 95\% interval. The product
$(\text{critical value})\times\SE$ is the \emph{margin of error}.
\end{definition}

Three ingredients change from case to case. The \emph{estimate} ($\xbar$,
$\hat p$, a difference of means). Its \emph{standard error} (always derived from
the sampling distribution). The \emph{critical value} ($1.96$ for a 95\% Normal
interval, a slightly larger $t$ value when $\sigma$ is estimated).

\section{The z-interval from a pivot}
\label{sec:ci-pivot}

The informal argument above can be made exact. Assume $X_1,\dots,X_n \iid
\Normal(\mu,\sigma^2)$ with $\sigma$ known (or $n$ large enough for the CLT).
Define
\[
  Z = \frac{\Xbar - \mu}{\sigma/\sqrt n}.
\]
$Z$ depends on the unknown $\mu$, yet its distribution, $\Normal(0,1)$, does
not depend on $\mu$ or anything else unknown. A quantity with this property is
called a \emph{pivot}. Because its distribution is fully known, we can write an
exact probability statement about it:
\[
  \P\left(-1.96 \le \frac{\Xbar-\mu}{\sigma/\sqrt n} \le 1.96\right) = 0.95 .
\]
Now rearrange the inequalities to isolate $\mu$. Multiply through by
$\sigma/\sqrt n$:
\[
  -1.96\frac{\sigma}{\sqrt n} \le \Xbar - \mu \le 1.96\frac{\sigma}{\sqrt n}.
\]
Subtract $\Xbar$ and multiply by $-1$ (which flips both inequalities):
\[
  \Xbar - 1.96\frac{\sigma}{\sqrt n} \le \mu \le \Xbar + 1.96\frac{\sigma}{\sqrt n}.
\]
The event has not changed, only how it is written. So its probability is still
0.95.

\begin{theorem}[z-interval for a mean]
If $\Xbar$ is Normal (exactly, or approximately by the CLT) and $\sigma$ is
known, then
\[
  \Xbar \pm z_{\alpha/2}\,\frac{\sigma}{\sqrt n}
\]
contains $\mu$ with probability $1-\alpha$, where $z_{\alpha/2}$ is the point
with upper-tail area $\alpha/2$ under $\Normal(0,1)$: $z_{0.05}=1.645$ (90\%),
$z_{0.025}=1.960$ (95\%), $z_{0.005}=2.576$ (99\%).
\end{theorem}

\begin{figure}[H]
\centering
\begin{tikzpicture}[x=0.95cm, y=1cm]
  % top: distance statement around mu
  \draw[->, thick] (-6,2) -- (6.3,2) node[right] {\small ms};
  \fill[histfill!50] (-1.96*1.4,1.85) rectangle (1.96*1.4,2.15);
  \draw[thick, distblue] (0,1.75) -- (0,2.25) node[above] {$\mu$};
  \node[font=\small, distblue, anchor=south] at (2.744,2.15) {$\mu+1.96\,\SE$};
  \node[font=\small, distblue, anchor=south] at (-2.744,2.15) {$\mu-1.96\,\SE$};
  \fill[quotientorange] (1.2,2) circle (2.5pt);
  \node[quotientorange, font=\small, anchor=north] at (1.2,1.8) {$\xbar$};
  \node[anchor=west, font=\small, align=left] at (-6,2.95) {$\xbar$ lands inside the band around $\mu$ (prob.\ 0.95)};
  % bottom: interval around xbar
  \draw[->, thick] (-6,0) -- (6.3,0) node[right] {\small ms};
  \draw[very thick, quotientorange] (1.2-2.744,0) -- (1.2+2.744,0);
  \draw[very thick, quotientorange] (1.2-2.744,-0.15) -- (1.2-2.744,0.15);
  \draw[very thick, quotientorange] (1.2+2.744,-0.15) -- (1.2+2.744,0.15);
  \fill[quotientorange] (1.2,0) circle (2.5pt);
  \node[quotientorange, font=\small, anchor=north] at (1.2,-0.2) {$\xbar$};
  \draw[thick, distblue] (0,-0.25) -- (0,0.25) node[above] {$\mu$};
  \node[anchor=west, font=\small] at (-6,-0.9) {$\iff$ the interval $\xbar \pm 1.96\,\SE$ covers $\mu$};
  \draw[dashed, gray] (0,0.8) -- (0,1.55);
\end{tikzpicture}
\caption{The pivot argument in one picture. ``$\xbar$ is within $1.96\,\SE$ of
$\mu$'' and ``$\mu$ is within $1.96\,\SE$ of $\xbar$'' are the same event, so
they have the same probability.}
\label{fig:ci-pivot}
\end{figure}

\begin{whybox}
Why $1.96$? Because $\P(\abs{Z}\le 1.96) = 0.95$ for a standard Normal. The
critical value is not a convention; it is whatever number traps the chosen
fraction of the pivot's distribution. A 99\% interval needs to trap 99\%, so
it reaches out to $2.576$. Nothing else about the construction changes.
\end{whybox}

\begin{recoverybox}
Pause and think: with the same 36 page loads, what is the 99\% interval?
(Answer: $212 \pm 2.576\times 5 = 212\pm 12.9 = [199.1,\ 224.9]$ ms. Wider,
because it must succeed more often.)
\end{recoverybox}

\section{What ``95\% confidence'' means}
\label{sec:ci-coverage}

The probability statement above is about $\Xbar$, a random variable, before
sampling. Once we have computed $[202.2, 221.8]$, nothing is random any more:
$\mu$ is a fixed number and it either is or is not in that range. The 95\%
describes the \emph{procedure}, not the particular interval.

The cleanest way to see this is to run the procedure many times. We simulated
a population with $\mu = 200$ ms and $\sigma=30$ ms, drew 25 independent
samples of $n=36$, and built the 95\% z-interval $\xbar \pm 9.8$ from each.

\begin{figure}[H]
\centering
\begin{tikzpicture}
\begin{axis}[
  width=12cm, height=9cm,
  xmin=176, xmax=240, ymin=0, ymax=27.5,
  xtick={180,190,200,210,220,230},
  xlabel={ms}, ylabel={sample number},
  ytick={1,5,10,15,20,25},
  axis lines=left,
  legend style={at={(0.98,0.02)}, anchor=south east, font=\small},
]
\draw[thick, dashed, black] (axis cs:200,0) -- (axis cs:200,26);
\addplot[only marks, mark=*, mark size=1.4pt, sourceblue,
  error bars/.cd, x dir=both, x fixed=9.8, error bar style={sourceblue, thick}]
coordinates {
(202.33,1) (192.17,2) (200.63,3) (207.00,4) (202.67,5) (203.06,6) (194.21,7)
(193.18,9) (203.30,10) (201.10,11) (198.45,13) (204.39,14) (204.04,15)
(199.14,16) (203.96,17) (204.22,18) (205.60,19) (195.73,20) (194.95,21)
(196.78,22) (200.97,23) (191.50,24) (200.69,25)
};
\addlegendentry{covers $\mu$ (23)}
\addplot[only marks, mark=*, mark size=1.6pt, trapred,
  error bars/.cd, x dir=both, x fixed=9.8, error bar style={trapred, very thick}]
coordinates { (211.60,8) (210.76,12) };
\addlegendentry{misses $\mu$ (2)}
\node[anchor=south, font=\small] at (axis cs:200,26) {$\mu=200$};
\end{axis}
\end{tikzpicture}
\caption{Twenty-five 95\% intervals from 25 simulated samples ($n=36$,
$\sigma=30$, true $\mu=200$). Each interval has the same width (19.6 ms); only
its centre moves. Samples 8 and 12 drew unusually slow requests and their
intervals (red) miss $\mu$. Over the long run, 5\% of intervals miss; here 2 of
25 did (8\%), which is within ordinary chance.}
\label{fig:ci-coverage}
\end{figure}

\begin{definition}[Coverage]
The \emph{coverage probability} of an interval procedure is the probability,
over repeated samples, that the interval contains the true parameter. A
procedure is a valid 95\% interval if its coverage is (at least
approximately) 0.95 for every possible value of the parameter.
\end{definition}

\begin{intuitionbox}
Think of a ring toss. The peg ($\mu$) is fixed. Each sample throws a ring of
fixed width at it. A 95\% procedure is a thrower who rings the peg 95\% of the
time. After a single throw you do not get to see the peg; you only know the
thrower's success rate.
\end{intuitionbox}

\begin{trapbox}
Statements that are \textbf{wrong} for the interval $[202.2, 221.8]$:
\begin{itemize}
  \item ``There is a 95\% probability that $\mu$ lies between 202.2 and
        221.8.'' In frequentist statistics $\mu$ is not random; the probability is 0 or 1.
        (A Bayesian \emph{credible} interval does make this statement; see Chapter~14.)
  \item ``95\% of page loads take between 202.2 and 221.8 ms.'' The interval is
        about the \emph{mean}. Individual loads vary with $\sigma=30$; a range
        covering 95\% of loads would be roughly $212\pm 59$.
  \item ``If we repeat the study, 95\% of new sample means will fall in
        $[202.2,221.8]$.'' New means scatter around $\mu$, not around this
        $\xbar$. That fraction is lower, about 83\% on average.
  \item ``The interval contains all plausible values with certainty.'' One
        in twenty intervals misses, and you never know which one.
\end{itemize}
The right statement: ``We used a method that captures the true mean in 95\%
of samples; applied to our sample, it gives $[202.2, 221.8]$.''
\end{trapbox}

\begin{recoverybox}
Pause and think: in Figure~\ref{fig:ci-coverage}, if we had drawn 1000
samples instead of 25, about how many intervals would be red? (About 50. The
count of misses is itself $\Bin(1000, 0.05)$, so anywhere from roughly
37 to 64 is typical.)
\end{recoverybox}

\subsection{Confidence level versus width}

Higher confidence is not free. With the same data, raising the level raises the
critical value, so the interval widens. Using $\xbar=212$ and $\SE=5$:

\begin{center}
\begin{tabular}{lccc}
\toprule
Level & $z_{\alpha/2}$ & Margin of error & Interval (ms) \\
\midrule
80\%   & 1.282 & 6.41  & $[205.59,\ 218.41]$ \\
90\%   & 1.645 & 8.22  & $[203.78,\ 220.22]$ \\
95\%   & 1.960 & 9.80  & $[202.20,\ 221.80]$ \\
99\%   & 2.576 & 12.88 & $[199.12,\ 224.88]$ \\
99.9\% & 3.291 & 16.45 & $[195.55,\ 228.45]$ \\
\bottomrule
\end{tabular}
\end{center}

\begin{figure}[H]
\centering
\begin{tikzpicture}
\begin{axis}[
  width=11cm, height=5.6cm,
  xmin=192, xmax=232, ymin=0.3, ymax=5.7,
  ytick={1,2,3,4,5}, yticklabels={80\%,90\%,95\%,99\%,99.9\%},
  xlabel={mean page-load time (ms)}, ylabel={confidence level},
  axis lines=left,
]
\addplot[only marks, mark=*, mark size=1.6pt, structurepurple,
  error bars/.cd, x dir=both, x explicit, error bar style={structurepurple, thick}]
coordinates {
(212,1) +- (6.41,0)
(212,2) +- (8.22,0)
(212,3) +- (9.80,0)
(212,4) +- (12.88,0)
(212,5) +- (16.45,0)
};
\draw[dashed, gray] (axis cs:205,0.3) -- (axis cs:205,5.7);
\node[gray, font=\small, anchor=north east] at (axis cs:205,5.6) {last week: 205};
\end{axis}
\end{tikzpicture}
\caption{Same data, five confidence levels. Going from 95\% to 99.9\%
widens the interval by 68\%. The only ways to get a narrow interval with high
confidence are more data or less noisy data.}
\label{fig:ci-level-width}
\end{figure}

A 100\% confidence interval exists: $(-\infty, \infty)$. It is useless. Every
choice of level is a trade between being right more often and saying something
specific. 95\% is a convention, not a law; 90\% is common in engineering
dashboards and 99\% where errors are expensive.

\section{When \texorpdfstring{$\sigma$}{sigma} is unknown: the t-distribution}
\label{sec:ci-t}

In practice we almost never know $\sigma$. The natural fix is to plug in the
sample standard deviation $s$ (Chapter~2):
\[
  T = \frac{\Xbar - \mu}{s/\sqrt n}.
\]
This is no longer standard Normal. The denominator is now random too. With a
small sample, $s$ is a noisy estimate of $\sigma$; when $s$ happens to come
out low, $T$ gets inflated. Those occasional inflations give $T$ heavier tails
than $Z$.

How much does this matter? If we ignore it and use $\xbar \pm 1.96\,s/\sqrt n$
on Normal data, the true coverage is:

\begin{figure}[H]
\centering
\begin{tikzpicture}
\begin{axis}[
  width=11cm, height=5.6cm,
  xmode=log, log basis x=10,
  xmin=1.8, xmax=120, ymin=0.68, ymax=0.98,
  xtick={2,3,5,10,20,30,50,100}, xticklabels={2,3,5,10,20,30,50,100},
  ytick={0.70,0.75,0.80,0.85,0.90,0.95},
  yticklabel style={/pgf/number format/fixed, /pgf/number format/precision=2},
  xlabel={sample size $n$}, ylabel={actual coverage},
  axis lines=left, grid=major, grid style={gray!20},
]
\addplot[densityred, thick, mark=*, mark size=1.5pt] coordinates {
(2,0.6997) (3,0.8109) (4,0.8551) (5,0.8784) (6,0.8927) (8,0.9092) (10,0.9184)
(15,0.9298) (20,0.9352) (30,0.9403) (50,0.9443) (100,0.9472)
};
\draw[dashed, targetgreen, thick] (axis cs:1.8,0.95) -- (axis cs:120,0.95);
\node[targetgreen, font=\small, anchor=south west] at (axis cs:2.1,0.95) {nominal 95\%};
\end{axis}
\end{tikzpicture}
\caption{Exact coverage of the naive interval $\xbar\pm 1.96\,s/\sqrt n$ for
Normal data. At $n=5$ it covers only 88\% of the time; at $n=2$, 70\%. The
shortfall is caused entirely by the noise in $s$.}
\label{fig:ci-naive-coverage}
\end{figure}

\begin{historybox}
In 1908 William Sealy Gosset, a chemist at the Guinness brewery in Dublin,
had to judge barley and hop quality from tiny batches: often four or five
measurements. The practice of the day was to plug $s$ into the Normal formula as if it were
$\sigma$, and with $n=4$ that gave badly overconfident intervals, exactly the shortfall in
Figure~\ref{fig:ci-naive-coverage}. Gosset worked out (partly by hand
simulation with numbers drawn from shuffled cards) the distribution of
$(\Xbar-\mu)/(s/\sqrt n)$. Guinness did not allow staff to publish under their
own names, so he signed the paper ``Student''. The result made honest
inference from small samples possible, and his ratio is still called Student's
$t$.
\end{historybox}

\begin{theorem}[Student's t]
If $X_1,\dots,X_n \iid \Normal(\mu,\sigma^2)$, then
\[
  T = \frac{\Xbar-\mu}{s/\sqrt n} \sim t_{n-1},
\]
the $t$-distribution with $n-1$ degrees of freedom. It is symmetric around 0,
bell-shaped, has heavier tails than $\Normal(0,1)$, and approaches
$\Normal(0,1)$ as $n\to\infty$.
\end{theorem}

The \emph{degrees of freedom} (df) is the single number that picks out which
member of the $t$ family you have: small df means a squat curve with fat tails,
large df means nearly Normal. Here df $=n-1$ for the same reason $s^2$ divides by $n-1$:
the $n$ deviations $x_i - \xbar$ sum to zero, so only $n-1$ of them carry
independent information about spread (Chapter~2). The proof that $T$ has
exactly this distribution uses the independence of $\Xbar$ and $s^2$ for Normal
data; see Probability From the Ground Up for the derivation.

\begin{figure}[H]
\centering
\begin{tikzpicture}
\begin{axis}[
  width=12cm, height=6cm,
  domain=-4.5:4.5, samples=141,
  axis lines=left, ymin=0, ymax=0.43,
  xlabel={$t$}, ylabel={density},
  legend style={at={(0.98,0.97)}, anchor=north east, font=\small},
]
\addplot[black, very thick] {0.39894*exp(-x^2/2)};
\addlegendentry{$\Normal(0,1)$}
\addplot[distblue, thick] {0.389108*(1+x^2/10)^(-5.5)};
\addlegendentry{$t_{10}$}
\addplot[quotientorange, thick, dashed] {0.367553*(1+x^2/3)^(-2)};
\addlegendentry{$t_{3}$}
\addplot[densityred, thick, dotted] {0.318310/(1+x^2)};
\addlegendentry{$t_{1}$}
\draw[->, gray] (axis cs:3.3,0.12) -- (axis cs:3.0,0.03);
\node[gray, font=\small, anchor=south] at (axis cs:3.3,0.12) {heavier tails};
\end{axis}
\end{tikzpicture}
\caption{The $t$-distribution for 1, 3 and 10 degrees of freedom against the
standard Normal. Fewer degrees of freedom means a lower peak and fatter tails.
At $t=3$, the two-sided tail area is 0.27\% for the Normal but 5.8\% for
$t_3$. By $t_{30}$ the curves are nearly indistinguishable.}
\label{fig:ci-t-shapes}
\end{figure}

Fatter tails mean we must reach further out to trap 95\%. The \emph{critical
value} (the cut-off on the horizontal axis that leaves the chosen tail area
outside) becomes $t_{\alpha/2,\,n-1}$ in place of $z_{\alpha/2}$:

\begin{center}
\begin{tabular}{rccc}
\toprule
df & 90\% ($t_{0.05}$) & 95\% ($t_{0.025}$) & 99\% ($t_{0.005}$) \\
\midrule
1   & 6.314 & 12.706 & 63.657 \\
2   & 2.920 & 4.303 & 9.925 \\
3   & 2.353 & 3.182 & 5.841 \\
4   & 2.132 & 2.776 & 4.604 \\
5   & 2.015 & 2.571 & 4.032 \\
6   & 1.943 & 2.447 & 3.707 \\
7   & 1.895 & 2.365 & 3.499 \\
8   & 1.860 & 2.306 & 3.355 \\
9   & 1.833 & 2.262 & 3.250 \\
10  & 1.812 & 2.228 & 3.169 \\
15  & 1.753 & 2.131 & 2.947 \\
20  & 1.725 & 2.086 & 2.845 \\
30  & 1.697 & 2.042 & 2.750 \\
50  & 1.676 & 2.009 & 2.678 \\
100 & 1.660 & 1.984 & 2.626 \\
$\infty$ & 1.645 & 1.960 & 2.576 \\
\bottomrule
\end{tabular}
\end{center}

The last row is the Normal. The table makes the practical point: below about
$n=30$ the difference matters; above it, $t$ and $z$ agree to within a few
percent.

\begin{theorem}[t-interval for a mean]
For a sample of size $n$ with mean $\xbar$ and standard deviation $s$,
\[
  \xbar \pm t_{\alpha/2,\,n-1}\,\frac{s}{\sqrt n}
\]
has coverage exactly $1-\alpha$ for Normal data, and approximately $1-\alpha$
for other data when $n$ is moderate and the data are not heavily skewed.
\end{theorem}

\subsection{Worked example}

\begin{walkbox}
Eight latency measurements (ms) from a staging server:
\[
  182,\ 195,\ 210,\ 177,\ 224,\ 201,\ 189,\ 214.
\]
\textbf{Step 1: mean.} The sum is 1592, so $\xbar = 1592/8 = 199$ ms.

\textbf{Step 2: standard deviation.} Deviations from 199:
$-17, -4, 11, -22, 25, 2, -10, 15$. Squares: $289, 16, 121, 484, 625, 4, 100,
225$, summing to 1864. So
\[
  s^2 = \frac{1864}{8-1} = 266.29, \qquad s = 16.32 \text{ ms}.
\]
\textbf{Step 3: standard error.} $\SE = 16.32/\sqrt 8 = 16.32/2.828 = 5.77$ ms.

\textbf{Step 4: critical value.} $\text{df} = 7$, so $t_{0.025,7} = 2.365$
from the table.

\textbf{Step 5: interval.} Margin $= 2.365 \times 5.77 = 13.64$ ms, giving
\[
  199 \pm 13.64 = [185.4,\ 212.6] \text{ ms}.
\]
Using $1.96$ instead would give $[187.7, 210.3]$, an interval 17\% too narrow,
with real coverage about 91\% instead of 95\%.
\end{walkbox}

\begin{lstlisting}[language=Python]
import numpy as np
from scipy import stats
x = np.array([182, 195, 210, 177, 224, 201, 189, 214])
n, m, s = len(x), x.mean(), x.std(ddof=1)     # ddof=1 -> divide by n-1
t = stats.t.ppf(0.975, df=n-1)                 # 2.365
print(m - t*s/np.sqrt(n), m + t*s/np.sqrt(n))  # 185.36 212.64
# one call: stats.t.interval(0.95, n-1, loc=m, scale=s/np.sqrt(n))
\end{lstlisting}

\begin{trapbox}
The t-interval assumes the \emph{sample mean} is roughly Normal, not the data.
For small $n$ with strong skew (latencies often have a long right tail), the
t-interval can still under-cover. Look at the data first (Chapter~3). For
skewed data, consider a log transform, an interval for the median, or the
bootstrap (Chapter~8).
\end{trapbox}

\section{Interval for a proportion}
\label{sec:ci-proportion}

A landing page gets $n=400$ visitors, of whom $x=46$ sign up. The sample
proportion is $\hat p = 46/400 = 0.115$. Each visitor is a Bernoulli trial, so
(Chapter~6)
\[
  \SE(\hat p) = \sqrt{\frac{p(1-p)}{n}}.
\]
We do not know $p$, so the simplest move is to plug in $\hat p$.

\begin{definition}[Wald interval]
\[
  \hat p \pm z_{\alpha/2}\sqrt{\frac{\hat p(1-\hat p)}{n}}.
\]
\end{definition}

For the landing page: $\SE = \sqrt{0.115\times 0.885/400} = 0.01595$, margin
$= 1.96\times 0.01595 = 0.0313$, interval $[0.084,\ 0.146]$. Somewhere between
8.4\% and 14.6\% of visitors sign up.

\subsection{Why Wald fails near 0 and 1}

Now a smaller test: 25 runs of a flaky integration test, 0 failures. Then
$\hat p = 0$, $\SE = \sqrt{0\times 1/25} = 0$, and the Wald interval is
$[0,0]$. It claims certainty that the failure rate is exactly zero after 25
runs. With 2 failures out of 25 it gives $0.08 \pm 0.106 = [-0.026,\ 0.186]$,
which includes impossible negative rates.

Two things go wrong at once:
\begin{enumerate}
  \item \textbf{The SE is estimated badly.} When $\hat p$ is near 0, the plug-in
        $\hat p(1-\hat p)$ is small, so the interval is too narrow exactly when
        $\hat p$ happens to have come out low.
  \item \textbf{The Normal shape is wrong.} A Binomial with small $np$ is
        strongly right-skewed and cannot go below zero. A symmetric interval
        cannot match it.
\end{enumerate}

We can compute the Wald interval's true coverage exactly, without simulation.
For fixed $n$ and $p$, sum the Binomial probabilities of every count $x$ whose
interval contains $p$. Figure~\ref{fig:ci-wald-wilson-coverage} does this for
$n=30$ and $p$ from 0.01 to 0.99.

\begin{figure}[H]
\centering
\begin{tikzpicture}
\begin{axis}[
  width=12.4cm, height=6.6cm,
  xmin=0, xmax=1, ymin=0.2, ymax=1.0,
  xlabel={true proportion $p$}, ylabel={coverage},
  ytick={0.2,0.4,0.6,0.8,0.95},
  axis lines=left,
  legend style={at={(0.5,0.03)}, anchor=south, font=\small, legend columns=2},
]
\addplot[densityred, thick] coordinates {
(0.01,0.260) (0.02,0.454) (0.03,0.597) (0.04,0.705) (0.05,0.782) (0.06,0.842) (0.07,0.883) (0.08,0.910) (0.09,0.937) (0.10,0.809) (0.11,0.853) (0.12,0.883) (0.13,0.904) (0.14,0.930) (0.15,0.942) (0.16,0.865) (0.17,0.897) (0.18,0.914) (0.19,0.925) (0.20,0.946) (0.21,0.888) (0.22,0.902) (0.23,0.927) (0.24,0.936) (0.25,0.941) (0.26,0.910) (0.27,0.921) (0.28,0.928) (0.29,0.949) (0.30,0.953) (0.31,0.916) (0.32,0.921) (0.33,0.945) (0.34,0.947) (0.35,0.911) (0.36,0.937) (0.37,0.941) (0.38,0.942) (0.39,0.929) (0.40,0.935) (0.41,0.938) (0.42,0.937) (0.43,0.930) (0.44,0.934) (0.45,0.935) (0.46,0.933) (0.47,0.931) (0.48,0.934) (0.49,0.933) (0.50,0.957) (0.51,0.933) (0.52,0.934) (0.53,0.931) (0.54,0.933) (0.55,0.935) (0.56,0.934) (0.57,0.930) (0.58,0.937) (0.59,0.938) (0.60,0.935) (0.61,0.929) (0.62,0.942) (0.63,0.941) (0.64,0.937) (0.65,0.911) (0.66,0.947) (0.67,0.945) (0.68,0.921) (0.69,0.916) (0.70,0.953) (0.71,0.949) (0.72,0.928) (0.73,0.921) (0.74,0.910) (0.75,0.941) (0.76,0.936) (0.77,0.927) (0.78,0.902) (0.79,0.888) (0.80,0.946) (0.81,0.925) (0.82,0.914) (0.83,0.897) (0.84,0.865) (0.85,0.942) (0.86,0.930) (0.87,0.904) (0.88,0.883) (0.89,0.853) (0.90,0.809) (0.91,0.937) (0.92,0.910) (0.93,0.883) (0.94,0.842) (0.95,0.782) (0.96,0.705) (0.97,0.597) (0.98,0.454) (0.99,0.260)
};
\addlegendentry{Wald}
\addplot[distblue, thick] coordinates {
(0.01,0.964) (0.02,0.978) (0.03,0.940) (0.04,0.969) (0.05,0.939) (0.06,0.968) (0.07,0.945) (0.08,0.971) (0.09,0.952) (0.10,0.974) (0.11,0.959) (0.12,0.956) (0.13,0.951) (0.14,0.939) (0.15,0.965) (0.16,0.954) (0.17,0.951) (0.18,0.948) (0.19,0.940) (0.20,0.964) (0.21,0.957) (0.22,0.954) (0.23,0.953) (0.24,0.948) (0.25,0.968) (0.26,0.942) (0.27,0.939) (0.28,0.961) (0.29,0.957) (0.30,0.930) (0.31,0.954) (0.32,0.952) (0.33,0.947) (0.34,0.947) (0.35,0.947) (0.36,0.943) (0.37,0.964) (0.38,0.942) (0.39,0.940) (0.40,0.962) (0.41,0.938) (0.42,0.937) (0.43,0.959) (0.44,0.958) (0.45,0.935) (0.46,0.957) (0.47,0.957) (0.48,0.934) (0.49,0.956) (0.50,0.957) (0.51,0.956) (0.52,0.934) (0.53,0.957) (0.54,0.957) (0.55,0.935) (0.56,0.958) (0.57,0.959) (0.58,0.937) (0.59,0.938) (0.60,0.962) (0.61,0.940) (0.62,0.942) (0.63,0.964) (0.64,0.943) (0.65,0.947) (0.66,0.947) (0.67,0.947) (0.68,0.952) (0.69,0.954) (0.70,0.930) (0.71,0.957) (0.72,0.961) (0.73,0.939) (0.74,0.942) (0.75,0.968) (0.76,0.948) (0.77,0.953) (0.78,0.954) (0.79,0.957) (0.80,0.964) (0.81,0.940) (0.82,0.948) (0.83,0.951) (0.84,0.954) (0.85,0.965) (0.86,0.939) (0.87,0.951) (0.88,0.956) (0.89,0.959) (0.90,0.974) (0.91,0.952) (0.92,0.971) (0.93,0.945) (0.94,0.968) (0.95,0.939) (0.96,0.969) (0.97,0.940) (0.98,0.978) (0.99,0.964)
};
\addlegendentry{Wilson}
\draw[dashed, black] (axis cs:0,0.95) -- (axis cs:1,0.95);
\end{axis}
\end{tikzpicture}
\caption{Exact coverage of nominal 95\% intervals for a proportion, $n=30$.
Wald (red) collapses near the edges: at $p=0.05$ it covers only 78\% of the
time, at $p=0.01$ only 26\%. Wilson (blue) stays between 93\% and 98\%
everywhere. Both curves are jagged because $x$ can only take 31 values.}
\label{fig:ci-wald-wilson-coverage}
\end{figure}

\subsection{The Wilson interval}

The fix is to stop plugging $\hat p$ into the standard error. Ask instead:
which values of $p$ are close enough to $\hat p$, measured in units of
\emph{their own} standard error?

\begin{definition}[Wilson score interval]
The set of $p$ satisfying
$\abs{\hat p - p} \le z\sqrt{p(1-p)/n}$ is the interval
\[
  \frac{\hat p + \dfrac{z^2}{2n}}{1+\dfrac{z^2}{n}}
  \;\pm\;
  \frac{z}{1+\dfrac{z^2}{n}}\sqrt{\frac{\hat p(1-\hat p)}{n} + \frac{z^2}{4n^2}},
\]
where $z = z_{\alpha/2}$.
\end{definition}

\begin{whybox}
Square both sides of $\abs{\hat p - p} \le z\sqrt{p(1-p)/n}$:
\[
  (\hat p - p)^2 \le \frac{z^2}{n}\,p(1-p).
\]
This is a quadratic inequality in $p$. Its two roots are the endpoints above.
The pivot $(\hat p - p)/\sqrt{p(1-p)/n}$ uses the \emph{true} $p$ in the SE,
so it is much closer to $\Normal(0,1)$ than the Wald version. Two
consequences: the centre is pulled from $\hat p$ toward $1/2$, and the interval
can never leave $[0,1]$.
\end{whybox}

For the flaky test, $x=0$, $n=25$: with $\hat p=0$ the formula simplifies to
$[0,\ (z^2/n)/(1+z^2/n)] = [0,\ 0.1537/1.1537] = [0,\ 0.133]$. Zero failures in
25 runs is consistent with a failure rate up to 13\%.

\begin{figure}[H]
\centering
\begin{tikzpicture}
\begin{axis}[
  width=11.5cm, height=6.4cm,
  xmin=-0.07, xmax=0.42, ymin=-0.6, ymax=5.6,
  ytick={0,1,2,3,4,5}, yticklabels={0/25,1/25,2/25,3/25,4/25,5/25},
  xtick={0,0.1,0.2,0.3,0.4},
  xticklabel style={/pgf/number format/fixed},
  xlabel={proportion}, ylabel={observed $x/n$},
  axis lines=left,
  legend style={at={(0.98,0.03)}, anchor=south east, font=\small},
]
\fill[trapred!12] (axis cs:-0.07,-0.6) rectangle (axis cs:0,5.6);
\draw[gray] (axis cs:0,-0.6) -- (axis cs:0,5.6);
\addplot[only marks, mark=*, mark size=1.4pt, densityred,
  error bars/.cd, x dir=both, x explicit, error bar style={densityred, thick}]
coordinates {
(0.00,0.15) +- (0,0)
(0.04,1.15) +- (0.0768,0)
(0.08,2.15) +- (0.1063,0)
(0.12,3.15) +- (0.1274,0)
(0.16,4.15) +- (0.1437,0)
(0.20,5.15) +- (0.1568,0)
};
\addlegendentry{Wald}
\addplot[only marks, mark=square*, mark size=1.4pt, distblue,
  error bars/.cd, x dir=both, x explicit, error bar style={distblue, thick}]
coordinates {
(0.0666,-0.15) +- (0.0666,0)
(0.1013,0.85) +- (0.0941,0)
(0.1359,1.85) +- (0.1138,0)
(0.1706,2.85) +- (0.1290,0)
(0.2053,3.85) +- (0.1413,0)
(0.2400,4.85) +- (0.1514,0)
};
\addlegendentry{Wilson}
\node[trapred, font=\scriptsize, rotate=90] at (axis cs:-0.035,4.4) {impossible};
\end{axis}
\end{tikzpicture}
\caption{Wald (red, circle at $\hat p$) and Wilson (blue, square at its
centre) 95\% intervals for $x$ successes out of $n=25$. Wald is a single point
at $x=0$ and spills below zero for $x=1,2,3$. Wilson is asymmetric, shifted
toward $1/2$, and stays inside $[0,1]$.}
\label{fig:ci-wald-wilson-intervals}
\end{figure}

\begin{lstlisting}[language=Python]
from statsmodels.stats.proportion import proportion_confint
proportion_confint(46, 400, alpha=0.05, method="normal")  # Wald  (0.084, 0.146)
proportion_confint(46, 400, alpha=0.05, method="wilson")  # Wilson(0.087, 0.150)
proportion_confint(0, 25, alpha=0.05, method="wilson")    # (0.0, 0.133)
\end{lstlisting}

\begin{trapbox}
The textbook rule ``Wald is fine if $n\hat p \ge 10$ and $n(1-\hat p)\ge
10$'' is a minimum, not a guarantee. Even at $p=0.5$ with $n=30$, Wald's
coverage oscillates between 93\% and 96\%. Wilson costs nothing extra. Use it
by default, especially for error rates, click-through rates and anything
measured in fractions of a percent.
\end{trapbox}

\section{Margin of error and sample-size planning}
\label{sec:ci-sample-size}

Turn the interval around. Before collecting data, decide how precise the
answer must be, then solve for $n$. For a mean, the margin of error is
$E = z\,\sigma/\sqrt n$. Solving for $n$:
\[
  n = \left(\frac{z\,\sigma}{E}\right)^2,
\]
rounded \emph{up}. You need a guess for $\sigma$: from a pilot study, past
data, or a conservative upper bound.

\begin{example}
We want the mean page-load time to within $\pm 5$ ms at 95\% confidence, and
historical data say $\sigma \approx 30$ ms. Then
$n = (1.96\times 30/5)^2 = 11.76^2 = 138.3$, so 139 requests. To halve the
margin to $\pm 2.5$ ms we need $(1.96\times 30/2.5)^2 = 553.2$, so 554
requests: four times as many.
\end{example}

\begin{intuitionbox}
Precision costs quadratically. Because the margin shrinks like $1/\sqrt n$,
each extra decimal place of precision (a factor of 10) costs 100 times the
data.
\end{intuitionbox}

\subsection{What ``$\pm 3\%$'' in a poll means}

For a proportion, $E = z\sqrt{p(1-p)/n}$. The product $p(1-p)$ is largest at
$p=1/2$, where it equals $1/4$. Planning with $p=1/2$ is therefore
conservative:
\[
  E_{\max} = z\sqrt{\frac{1}{4n}} = \frac{z}{2\sqrt n}, \qquad
  n = \left(\frac{z}{2E}\right)^2.
\]
For $E = 0.03$ at 95\%: $n = (1.96/0.06)^2 = 1067.1$, so 1068 respondents. A
poll of 1000 people has $E_{\max} = 1.96/(2\sqrt{1000}) = 0.031$. That is
where the ubiquitous ``margin of error $\pm 3$ percentage points'' comes from.

\begin{figure}[H]
\centering
\begin{tikzpicture}
\begin{axis}[
  width=11.5cm, height=5.8cm,
  domain=100:4500, samples=120,
  xmin=0, xmax=4500, ymin=0, ymax=10.5,
  xlabel={sample size $n$}, ylabel={max.\ margin of error (\%)},
  axis lines=left, grid=major, grid style={gray!20},
  scaled x ticks=false,
  xticklabel style={/pgf/number format/fixed, /pgf/number format/1000 sep={,}},
]
\addplot[distblue, thick] {98/sqrt(x)};
\addplot[only marks, mark=*, trapred] coordinates {(1067,3.0) (100,9.8) (4268,1.5)};
\node[anchor=south west, font=\small, trapred] at (axis cs:1067,3.0) {$n=1067$: $\pm3.0\%$};
\node[anchor=west, font=\small, trapred] at (axis cs:130,9.8) {$n=100$: $\pm 9.8\%$};
\node[anchor=south east, font=\small, trapred] at (axis cs:4400,2.3) {$n=4268$: $\pm1.5\%$};
\node[anchor=west, font=\small, distblue] at (axis cs:1900,0.7) {$E=1.96/(2\sqrt n)$};
\end{axis}
\end{tikzpicture}
\caption{Worst-case 95\% margin of error for a proportion versus sample size.
The curve flattens fast: going from 100 to 1067 respondents cuts the margin
from 9.8 to 3 points, but halving it again to 1.5 points needs four times
as many, about 4300.}
\label{fig:ci-moe-curve}
\end{figure}

\begin{trapbox}
The poll margin covers \emph{only} random sampling error. It says nothing about
who refused to answer, leading question wording, or people lying (Chapter~5).
Those biases do not shrink with $n$ and are often larger than 3 points.
Also: in a two-way race the margin on the \emph{gap} between the candidates is
roughly twice the margin on each share, because the shares move in opposite
directions: every respondent who switches adds to one and subtracts from the
other.
\end{trapbox}

The population size does not appear in any of these formulas. A random sample
of 1000 gives the same precision for a city of 50{,}000 as for a country of
300 million, as long as the population is much larger than the sample.

\section{Comparing two groups}
\label{sec:ci-two-sample}

Most practical questions are comparisons: is the new build faster, does the
new button convert better? The estimate is a difference, and we need the SE
of a difference.

\begin{whybox}
For independent estimates $A$ and $B$,
$\Var(A - B) = \Var(A) + \Var(B)$. Variances add even when you subtract,
because a minus sign flips direction but not spread. So
\[
  \SE(A - B) = \sqrt{\SE(A)^2 + \SE(B)^2}.
\]
Standard errors combine like the sides of a right triangle, not by plain
addition.
\end{whybox}

\subsection{Difference in means (Welch)}

\begin{walkbox}
Old build: $n_A = 40$ requests, $\xbar_A = 212$ ms, $s_A = 30$ ms.
New build: $n_B = 35$ requests, $\xbar_B = 198$ ms, $s_B = 22$ ms.

\textbf{Estimate.} $\xbar_A - \xbar_B = 14$ ms.

\textbf{SE.} $s_A^2/n_A = 900/40 = 22.50$ and $s_B^2/n_B = 484/35 = 13.83$, so
$\SE = \sqrt{22.50 + 13.83} = \sqrt{36.33} = 6.03$ ms.

\textbf{Degrees of freedom} (Welch--Satterthwaite):
\[
  \nu = \frac{\left(\dfrac{s_A^2}{n_A}+\dfrac{s_B^2}{n_B}\right)^2}
             {\dfrac{(s_A^2/n_A)^2}{n_A-1}+\dfrac{(s_B^2/n_B)^2}{n_B-1}}
      = \frac{36.33^2}{\dfrac{22.50^2}{39}+\dfrac{13.83^2}{34}}
      = \frac{1319.9}{12.98 + 5.62} = 70.9 .
\]
\textbf{Interval.} $t_{0.025,\,70.9} = 1.994$, so
$14 \pm 1.994\times 6.03 = 14 \pm 12.0 = [2.0,\ 26.0]$ ms.

The new build is faster by somewhere between 2 and 26 ms. The interval excludes
zero, but only just: the data are compatible with a negligible speedup or a
large one.
\end{walkbox}

\begin{definition}[Welch interval]
\[
  (\xbar_A - \xbar_B) \pm t_{\alpha/2,\,\nu}\sqrt{\frac{s_A^2}{n_A} + \frac{s_B^2}{n_B}},
\]
with $\nu$ from the Welch--Satterthwaite formula above. It does not assume
the two groups have equal variances.
\end{definition}

The $\nu$ formula is an approximation: it picks the $t$ distribution whose
degrees of freedom best match how noisy the estimated SE is. You do not need to compute it by
hand; software does. It always lies between $\min(n_A,n_B)-1$ and
$n_A+n_B-2$, and using $\min(n_A,n_B)-1$ is a safe hand shortcut. Older texts
use a ``pooled'' interval that assumes equal variances. Welch loses almost
nothing when variances are equal and is far safer when they are not, so use
Welch by default.

\begin{lstlisting}[language=Python]
from scipy import stats
res = stats.ttest_ind(a, b, equal_var=False)   # Welch
print(res.confidence_interval(0.95))           # CI for mean(a) - mean(b)
\end{lstlisting}

\subsection{Difference in proportions}

An A/B test shows the old checkout page to 2400 users (120 buy, $\hat p_A =
0.050$) and the new page to 2400 users (156 buy, $\hat p_B = 0.065$). The
difference is $\hat p_B - \hat p_A = 0.015$, or 1.5 percentage points.
\[
  \SE = \sqrt{\frac{0.05\times0.95}{2400} + \frac{0.065\times0.935}{2400}}
      = \sqrt{0.00001979 + 0.00002532} = 0.00672.
\]
The 95\% interval is $0.015 \pm 1.96\times 0.00672 = 0.015\pm 0.0132 =
[0.0018,\ 0.0282]$. The new page lifts conversion by between 0.2 and 2.8
points. In relative terms, that is anywhere from a 4\% to a 56\% improvement on the
5\% baseline: the direction is fairly clear, the size is not.

\begin{figure}[H]
\centering
\begin{tikzpicture}
\begin{axis}[
  name=left, width=6.3cm, height=5.5cm,
  xmin=0.5, xmax=2.5, ymin=3.5, ymax=8,
  xtick={1,2}, xticklabels={old (A), new (B)},
  ylabel={conversion rate (\%)},
  axis lines=left, title style={font=\small}, title={ each group},
]
\addplot[only marks, mark=*, sourceblue,
  error bars/.cd, y dir=both, y explicit, error bar style={sourceblue, thick}]
coordinates { (1,5.0) +- (0,0.872) (2,6.5) +- (0,0.986) };
\end{axis}
\begin{axis}[
  at={(left.east)}, anchor=west, xshift=1.6cm,
  width=5.2cm, height=5.5cm,
  xmin=0.5, xmax=1.5, ymin=-1, ymax=3.5,
  xtick={1}, xticklabels={B $-$ A},
  ylabel={difference (points)},
  axis lines=left, title style={font=\small}, title={ difference},
]
\draw[dashed, gray] (axis cs:0.5,0) -- (axis cs:1.5,0);
\addplot[only marks, mark=diamond*, mark size=3pt, targetgreen,
  error bars/.cd, y dir=both, y explicit, error bar style={targetgreen, thick}]
coordinates { (1,1.5) +- (0,1.32) };
\end{axis}
\end{tikzpicture}
\caption{The A/B checkout test with 95\% intervals. Left: each page's
conversion rate. Right: the interval for the difference, $[0.18, 2.82]$
points, which excludes zero. The question ``is B better?'' is answered by
the right panel, not by squinting at the left.}
\label{fig:ci-ab-test}
\end{figure}

\subsection{Paired data}

Eight servers had their p50 latency measured before and after a kernel
upgrade (ms):

\begin{center}
\begin{tabular}{lcccccccc}
\toprule
Server & 1 & 2 & 3 & 4 & 5 & 6 & 7 & 8 \\
\midrule
Before & 241 & 198 & 225 & 260 & 210 & 233 & 247 & 219 \\
After  & 228 & 195 & 214 & 249 & 207 & 220 & 238 & 213 \\
Difference $d$ & 13 & 3 & 11 & 11 & 3 & 13 & 9 & 6 \\
\bottomrule
\end{tabular}
\end{center}

\begin{recoverybox}
Pause and think before reading on: should we compare the mean of the
``before'' row with the mean of the ``after'' row as two separate groups?
\end{recoverybox}

The two columns are not independent samples. Server 4 is slow both times;
server 2 is fast both times. The before and after columns have correlation
0.99. The right analysis works on the differences, which turns the problem into
a one-sample t-interval:
\[
  \bar d = 8.625, \quad s_d = 4.14, \quad \SE = 4.14/\sqrt 8 = 1.463, \quad
  t_{0.025,7} = 2.365,
\]
\[
  8.625 \pm 2.365\times 1.463 = 8.625 \pm 3.46 = [5.17,\ 12.08] \text{ ms}.
\]

If we wrongly treat the columns as independent groups (Welch), the SE becomes
$9.44$ ms and the interval is $[-11.7,\ 28.9]$ ms, which includes zero and is
six times wider. The server-to-server variation (SD about 20 ms) swamps the
upgrade effect (about 9 ms). Pairing subtracts it out.

\begin{figure}[H]
\centering
\begin{tikzpicture}
\begin{axis}[
  width=8cm, height=6.2cm,
  xmin=0.6, xmax=2.4, ymin=190, ymax=265,
  xtick={1,2}, xticklabels={before, after},
  ylabel={p50 latency (ms)},
  axis lines=left,
]
\addplot[structurepurple, mark=*, mark size=1.5pt] coordinates {(1,241) (2,228)};
\addplot[structurepurple, mark=*, mark size=1.5pt] coordinates {(1,198) (2,195)};
\addplot[structurepurple, mark=*, mark size=1.5pt] coordinates {(1,225) (2,214)};
\addplot[structurepurple, mark=*, mark size=1.5pt] coordinates {(1,260) (2,249)};
\addplot[structurepurple, mark=*, mark size=1.5pt] coordinates {(1,210) (2,207)};
\addplot[structurepurple, mark=*, mark size=1.5pt] coordinates {(1,233) (2,220)};
\addplot[structurepurple, mark=*, mark size=1.5pt] coordinates {(1,247) (2,238)};
\addplot[structurepurple, mark=*, mark size=1.5pt] coordinates {(1,219) (2,213)};
\end{axis}
\end{tikzpicture}
\caption{Before/after latency for eight servers, one line per server. The
spread \emph{between} servers (about 65 ms top to bottom) is large, yet every
line slopes down. A paired interval looks only at the slopes.}
\label{fig:ci-paired}
\end{figure}

\begin{trapbox}
Whenever each observation in group A has a natural partner in group B (same
user, same server, same patient, same day), use the paired interval. Using a
two-sample interval on paired data does not give a wrong centre, but it
throws away the pairing and can make a clear effect look like noise.
\end{trapbox}

\section{One-sided bounds}
\label{sec:ci-one-sided}

Sometimes only one direction matters. An SLO says mean latency must be below
220 ms; we want an \emph{upper confidence bound}: a number we are 95\%
confident the mean does not exceed. Put the whole $\alpha=5\%$ in one tail:
\[
  \text{upper bound} = \xbar + t_{\alpha,\,n-1}\frac{s}{\sqrt n}.
\]
For the eight staging latencies ($\xbar = 199$, $\SE=5.77$, df $=7$):
$t_{0.05,7} = 1.895$, so the bound is $199 + 1.895\times 5.77 = 209.9$ ms.
The SLO of 220 ms is met with 95\% confidence.

\begin{figure}[H]
\centering
\begin{tikzpicture}
\begin{axis}[
  name=two, width=6.4cm, height=4.6cm,
  domain=-4:4, samples=101, axis lines=left, ymin=0, ymax=0.45,
  ytick=\empty, xtick={-2.365,0,2.365}, xticklabels={$-2.365$,0,$2.365$},
  xticklabel style={font=\scriptsize},
  title style={font=\small}, title={ two-sided 95\%: $2.5\%$ each tail},
]
\addplot[fill=trapred!40, draw=none, domain=-4:-2.365] {0.38499*(1+x^2/7)^(-4)} \closedcycle;
\addplot[fill=trapred!40, draw=none, domain=2.365:4] {0.38499*(1+x^2/7)^(-4)} \closedcycle;
\addplot[distblue, thick] {0.38499*(1+x^2/7)^(-4)};
\end{axis}
\begin{axis}[
  at={(two.east)}, anchor=west, xshift=0.6cm,
  width=6.4cm, height=4.6cm,
  domain=-4:4, samples=101, axis lines=left, ymin=0, ymax=0.45,
  ytick=\empty, xtick={0,1.895}, xticklabels={0,$1.895$},
  xticklabel style={font=\scriptsize},
  title style={font=\small}, title={ one-sided 95\%: $5\%$ upper tail},
]
\addplot[fill=trapred!40, draw=none, domain=1.895:4] {0.38499*(1+x^2/7)^(-4)} \closedcycle;
\addplot[distblue, thick] {0.38499*(1+x^2/7)^(-4)};
\end{axis}
\end{tikzpicture}
\caption{The $t_7$ density. A two-sided 95\% interval splits the 5\% miss
rate between both tails and reaches out to 2.365. A one-sided bound puts all
5\% in one tail, so its critical value is smaller (1.895) and the bound is
tighter on the side you care about.}
\label{fig:ci-one-sided}
\end{figure}

\begin{trapbox}
Decide on a one-sided bound \emph{before} looking at the data. Picking the
side after seeing which way the estimate went is a way of quietly doubling
your error rate. The one-sided 95\% upper bound equals the upper end of a
two-sided 90\% interval, which shows exactly what you are giving up.
\end{trapbox}

\section{Error bars: SD, SE, or CI?}
\label{sec:ci-error-bars}

A bar chart with whiskers is meaningless until you know what the whiskers are.
Three common choices look alike and mean very different things. Take two
groups of $n=10$ task-completion times (seconds):
\begin{center}
\begin{tabular}{lcccc}
\toprule
Group & mean & SD & SE $=$ SD$/\sqrt{10}$ & 95\% CI half-width ($t_{0.025,9}=2.262$) \\
\midrule
A & 50.0 & 9.79 & 3.09 & 7.00 \\
B & 62.0 & 8.88 & 2.81 & 6.35 \\
\bottomrule
\end{tabular}
\end{center}

\begin{figure}[H]
\centering
\begin{tikzpicture}
\begin{axis}[
  width=12.2cm, height=6.4cm,
  xmin=0.2, xmax=8.8, ymin=30, ymax=92,
  xtick={1,2,4,5,7,8}, xticklabels={A,B,A,B,A,B},
  ylabel={completion time (s)},
  axis lines=left, clip=false,
]
\addplot[only marks, mark=o, mark size=1.2pt, residgray] coordinates {
(0.75,38) (0.80,52) (0.85,61) (0.75,45) (0.80,49) (0.85,70) (0.75,44) (0.80,55) (0.85,41) (0.70,45)
(1.75,58) (1.80,66) (1.85,49) (1.75,72) (1.80,61) (1.85,55) (1.75,80) (1.80,57) (1.85,63) (1.70,59)
};
\addplot[only marks, mark=*, sourceblue,
  error bars/.cd, y dir=both, y explicit, error bar style={sourceblue, thick}]
coordinates { (1.1,50) +- (0,9.79) (2.1,62) +- (0,8.88) };
\addplot[only marks, mark=*, quotientorange,
  error bars/.cd, y dir=both, y explicit, error bar style={quotientorange, thick}]
coordinates { (4,50) +- (0,3.09) (5,62) +- (0,2.81) };
\addplot[only marks, mark=*, targetgreen,
  error bars/.cd, y dir=both, y explicit, error bar style={targetgreen, thick}]
coordinates { (7,50) +- (0,7.00) (8,62) +- (0,6.35) };
\node[sourceblue, font=\small, align=center] at (axis cs:1.5,88) {$\pm 1$ SD\\(spread of data)};
\node[quotientorange, font=\small, align=center] at (axis cs:4.5,88) {$\pm 1$ SE\\(68\% CI, roughly)};
\node[targetgreen, font=\small, align=center] at (axis cs:7.5,88) {95\% CI\\($\pm 2.26$ SE)};
\draw[gray!50] (axis cs:3,30) -- (axis cs:3,82);
\draw[gray!50] (axis cs:6,30) -- (axis cs:6,82);
\end{axis}
\end{tikzpicture}
\caption{The same two groups with three kinds of error bars. SD bars (with the
raw points in grey) describe individuals and do not shrink with $n$. SE bars are
the shortest and look the most impressive, but they are only about a 68\%
interval. 95\% CI bars are the ones that support inference about the means.}
\label{fig:ci-error-bars}
\end{figure}

\begin{itemize}
  \item \textbf{SD bars} answer ``how variable are individuals?'' They are
        descriptive. Use them when the reader cares about the spread of
        outcomes (how long will \emph{my} task take?).
  \item \textbf{SE bars} answer ``how precisely is each mean known?'' in units
        of one SE. They are about $1/\sqrt n$ as long as SD bars, so they look
        tight. A $\pm 1$ SE bar covers the true mean only about 68\% of the time.
  \item \textbf{95\% CI bars} answer the same question at a stated level. They
        are what most readers think they are seeing.
\end{itemize}

\begin{trapbox}
Always label error bars in the caption: ``mean $\pm$ SD'', ``mean $\pm$ SE'',
or ``mean with 95\% CI''. An unlabelled error bar is uninterpretable, and
switching SD to SE is a common way (deliberate or not) to make noisy results
look solid.
\end{trapbox}

\subsection{The overlapping-intervals fallacy}

A tempting shortcut: ``the two 95\% intervals overlap, so the difference is not
significant.'' It is wrong.

Two services have mean latencies 200 ms and 207 ms, each with $\SE = 2$ ms.
Their 95\% intervals are $[196.1,\ 203.9]$ and $[203.1,\ 210.9]$. They overlap.
But the difference is $7$ ms with
\[
  \SE(\text{diff}) = \sqrt{2^2 + 2^2} = 2.83, \qquad
  7 \pm 1.96\times 2.83 = [1.5,\ 12.5],
\]
which excludes zero.

\begin{figure}[H]
\centering
\begin{tikzpicture}
\begin{axis}[
  name=grp, width=6.2cm, height=5.6cm,
  xmin=0.5, xmax=2.5, ymin=193, ymax=214,
  xtick={1,2}, xticklabels={service 1, service 2},
  ylabel={mean latency (ms)}, axis lines=left,
  title style={font=\small}, title={ 95\% CI of each mean},
]
\fill[quotientorange!25] (axis cs:0.5,203.08) rectangle (axis cs:2.5,203.92);
\addplot[only marks, mark=*, sourceblue,
  error bars/.cd, y dir=both, y explicit, error bar style={sourceblue, thick}]
coordinates { (1,200) +- (0,3.92) (2,207) +- (0,3.92) };
\node[quotientorange!80!black, font=\scriptsize, anchor=south] at (axis cs:1.5,204) {overlap};
\end{axis}
\begin{axis}[
  at={(grp.east)}, anchor=west, xshift=1.5cm,
  width=5cm, height=5.6cm,
  xmin=0.5, xmax=1.5, ymin=-2, ymax=14,
  xtick={1}, xticklabels={2 $-$ 1},
  ylabel={difference (ms)}, axis lines=left,
  title style={font=\small}, title={ 95\% CI of difference},
]
\draw[dashed, gray] (axis cs:0.5,0) -- (axis cs:1.5,0);
\addplot[only marks, mark=diamond*, mark size=3pt, targetgreen,
  error bars/.cd, y dir=both, y explicit, error bar style={targetgreen, thick}]
coordinates { (1,7) +- (0,5.54) };
\end{axis}
\end{tikzpicture}
\caption{Overlapping individual intervals (left) can coexist with a
difference interval that excludes zero (right). The SE of a difference is
$\sqrt 2\approx 1.41$ times one SE, not 2 times, so judging by overlap is too
conservative.}
\label{fig:ci-overlap}
\end{figure}

\begin{whybox}
Non-overlap of two 95\% intervals requires the gap to exceed
$1.96(\SE_1 + \SE_2)$. Significance of the difference requires the gap to exceed
$1.96\sqrt{\SE_1^2 + \SE_2^2}$, which is smaller (the hypotenuse is shorter
than the two legs combined). With equal SEs the ratio is $\sqrt 2$. So:
\begin{itemize}
  \item intervals do not overlap $\Rightarrow$ the difference is significant
        (with equal SEs, at about the 0.6\% level, not 5\%);
  \item intervals overlap $\Rightarrow$ you cannot conclude anything. Compute
        the interval for the difference.
\end{itemize}
For paired data, overlap is even less informative, because the individual
intervals ignore the pairing entirely.
\end{whybox}

\section{Choosing an interval}
\label{sec:ci-choosing}

\begin{figure}[H]
\centering
\begin{tikzpicture}[
  node distance=0.55cm and 0.5cm,
  q/.style={draw, rounded corners, fill=structurepurple!10, font=\small, align=center, text width=2.6cm},
  a/.style={draw, fill=targetgreen!12, font=\small, align=center, text width=2.3cm},
  >={Stealth[length=2mm]}
]
\node[q] (start) {What is the parameter?};
\node[q, below left=0.7cm and 0.9cm of start] (mean) {One mean or two?};
\node[q, below right=0.7cm and 0.9cm of start] (prop) {One proportion\\or two?};
\node[a, below left=0.7cm and -0.6cm of mean] (one) {$\xbar \pm t\,s/\sqrt n$\\(t-interval)};
\node[q, below right=0.7cm and -0.6cm of mean] (pair) {Paired?};
\node[a, below left=0.7cm and -1.2cm of pair] (paired) {t-interval on differences};
\node[a, below right=0.7cm and -1.2cm of pair] (welch) {Welch interval};
\node[a, below left=0.7cm and -0.6cm of prop] (wil) {Wilson interval};
\node[a, below right=0.7cm and -0.6cm of prop] (dp) {$\hat p_B - \hat p_A \pm z\,\SE$};
\draw[->] (start) -- node[above left, font=\scriptsize] {mean} (mean);
\draw[->] (start) -- node[above right, font=\scriptsize] {proportion} (prop);
\draw[->] (mean) -- node[left, font=\scriptsize] {one} (one);
\draw[->] (mean) -- node[right, font=\scriptsize] {two} (pair);
\draw[->] (pair) -- node[left, font=\scriptsize] {yes} (paired);
\draw[->] (pair) -- node[right, font=\scriptsize] {no} (welch);
\draw[->] (prop) -- node[left, font=\scriptsize] {one} (wil);
\draw[->] (prop) -- node[right, font=\scriptsize] {two} (dp);
\end{tikzpicture}
\caption{Choosing among the intervals of this chapter. For medians,
percentiles such as p99, and other statistics without a simple SE formula,
use the bootstrap (Chapter~8).}
\label{fig:ci-decision}
\end{figure}

\begin{connectionbox}
Machine-learning papers report ``accuracy $91.2 \pm 0.4$''. Check what the
$\pm$ is: often it is the SD across 5 random seeds, which measures training
instability, not the uncertainty in the true accuracy. A test-set accuracy is a
proportion, so a Wilson interval on (correct, test-set size) gives the
sampling uncertainty from the finite test set: with 2000 test examples and
91.2\% accuracy, the 95\% margin is about $\pm 1.2$ points, three times the
quoted $\pm 0.4$. Chapter~15 returns to this.
\end{connectionbox}

\section{Summary}
\label{sec:ci-summary}

\begin{center}
\renewcommand{\arraystretch}{1.35}
\begin{tabular}{lll}
\toprule
Parameter & Interval & Notes \\
\midrule
Mean, $\sigma$ known & $\xbar \pm z_{\alpha/2}\,\sigma/\sqrt n$ & rare in practice \\
Mean & $\xbar \pm t_{\alpha/2,n-1}\,s/\sqrt n$ & default for means \\
Proportion & Wilson & avoid Wald near 0 or 1 \\
Two means & $\Delta\xbar \pm t_{\alpha/2,\nu}\sqrt{s_A^2/n_A + s_B^2/n_B}$ & Welch $\nu$ \\
Two proportions & $\Delta\hat p \pm z_{\alpha/2}\sqrt{\sum \hat p(1-\hat p)/n}$ & A/B tests \\
Paired & $\bar d \pm t_{\alpha/2,n-1}\,s_d/\sqrt n$ & use differences \\
One-sided upper & $\xbar + t_{\alpha,n-1}\,s/\sqrt n$ & all of $\alpha$ in one tail \\
\bottomrule
\end{tabular}
\end{center}

\begin{itemize}
  \item Every interval is estimate $\pm$ critical value $\times$ SE.
  \item ``95\% confidence'' is the long-run hit rate of the method. A computed
        interval either contains $\mu$ or does not.
  \item Use $t$ instead of $z$ whenever $s$ replaces $\sigma$; it matters for
        $n\lesssim 30$.
  \item Higher confidence means wider intervals; precision improves only like
        $\sqrt n$. Plan with $n = (z\sigma/E)^2$ or $n=(z/2E)^2$ for
        proportions.
  \item SEs of independent estimates combine like the sides of a right
        triangle: $\sqrt{\SE_A^2+\SE_B^2}$. Pair when you can.
  \item Label error bars. Overlapping CIs do not imply ``no difference''.
\end{itemize}

Every interval here needed a formula for the SE. For a median, a p99, or a
ratio of two means there is no tidy formula. Chapter~8 gets the SE, and the
interval, by resampling the data itself: the bootstrap.

\section{Exercises}
\label{sec:ci-exercises}

\begin{exercise}
A sample of $n=64$ has $\xbar = 50$ and the population SD is known to be
$\sigma = 8$. Give the 95\% z-interval for $\mu$.
\end{exercise}

\begin{exercise}
An analyst quadruples the sample size. By what factor does the width of the
95\% interval change? What about if she instead raises the level from 95\% to
99\% at the original $n$?
\end{exercise}

\begin{exercise}
Six battery lifetimes (hours): $4.1, 3.8, 5.0, 4.6, 4.4, 3.9$. Compute a 95\%
t-interval for the mean lifetime.
\end{exercise}

\begin{exercise}
In a survey of 200 users, 30 report a bug. Compute the 95\% Wald interval for
the proportion of users who hit the bug.
\end{exercise}

\begin{exercise}
A deployment pipeline ran 50 times with no failures. What does the Wald
interval say? What is the 95\% Wilson upper bound on the failure rate?
\end{exercise}

\begin{exercise}
You want to estimate mean session length to within $\pm 3$ minutes with 99\%
confidence; a pilot suggests $\sigma \approx 15$ minutes. How many sessions do
you need?
\end{exercise}

\begin{exercise}
A poll of $n=1067$ reports $\pm 3\%$. Roughly how many respondents are needed for
$\pm 1.5\%$?
\end{exercise}

\begin{exercise}
True or false: (a) A 95\% CI of $[4.1, 4.8]$ means $\P(4.1\le\mu\le4.8)=0.95$.
(b) If we built 95\% intervals for 100 independent studies, we would expect
about 5 to miss their true parameter. (c) A 99\% interval from the same data is
always wider than the 95\% interval.
\end{exercise}

\begin{exercise}
Two group means each have a 95\% CI of half-width 4, and the means differ by 6.
The intervals overlap. Assuming independent groups with equal SEs, is the
difference significant at the 5\% level?
\end{exercise}

\subsection*{Answers}

\begin{enumerate}[label=\textbf{\arabic*.}]
  \item $\SE = 8/8 = 1$; $50 \pm 1.96 = [48.04,\ 51.96]$.
  \item Width scales as $1/\sqrt n$, so it halves. Going to 99\% multiplies
        the width by $2.576/1.960 = 1.31$.
  \item $\xbar = 4.30$, $s = 0.456$, $\SE = 0.186$, $t_{0.025,5}=2.571$;
        $4.30 \pm 0.479 = [3.82,\ 4.78]$ hours.
  \item $\hat p = 0.15$, $\SE = \sqrt{0.15\times0.85/200} = 0.0252$;
        $0.15 \pm 0.049 = [0.101,\ 0.199]$.
  \item Wald gives $[0,0]$, which is absurd. Wilson with $x=0$:
        upper $= (z^2/n)/(1+z^2/n) = 0.0768/1.0768 = 0.071$. The failure rate
        could plausibly be as high as 7\%.
  \item $n = (2.576\times 15/3)^2 = 12.88^2 = 165.9$, so 166 sessions.
  \item Halving the margin needs four times the sample: about $4\times 1067
        \approx 4270$.
  \item (a) False: $\mu$ is fixed; 95\% is the method's hit rate.
        (b) True. (c) True: $z_{0.005} > z_{0.025}$ (and likewise for $t$) with the same SE.
  \item Each SE is $4/1.96 = 2.04$; $\SE(\text{diff}) = 2.04\sqrt2 = 2.89$;
        the difference interval is $6 \pm 1.96\times 2.89 = 6\pm 5.66 =
        [0.34,\ 11.66]$, which excludes 0. Yes, significant, despite the
        overlap.
\end{enumerate}
